\uuid{trhY}
\chapitre{Statistique}
\niveau{L2}
\module{Analyse de données}
\sousChapitre{Tests d'hypothèses, intervalle de confiance}
\titre{Selon le choix de l'hypothèse}
\theme{tests statistiques}
\auteur{}
\datecreate{2022-09-28}
\organisation{AMSCC}
\difficulte{2}
\contenu{
\texte{$\mathcal N(\mu,\sigma)$ désigne la loi normale de moyenne $\mu$ et d’écart-type $\sigma$. Les quantiles sont à gauche : $t_{\nu;\gamma}$ pour la loi $St(\nu)$, avec une probabilité $\gamma$ à gauche du quantile.}
\texte{On compare les salaires de dix femmes de même ancienneté (3 à 5 ans) à une référence masculine fixée à 28 milliers d'euros. Les salaires observés, en milliers d'euros, sont : $24,27,31,19,26,27,22,15,33,21$.
On suppose les observations indépendantes de loi $\mathcal N(\mu,\sigma)$ ; l'incertitude de la référence masculine n'est pas étudiée ici.
M. A souhaite rechercher toute différence à 28 ; Mme B souhaite rechercher un salaire moyen inférieur à 28. On examine les deux démarches, supposées choisies avant l'observation des données.}
\begin{enumerate}
\item \question{Quelles hypothèses de modèle permettent un test exact de la moyenne ?}
\reponse{Un échantillon indépendant issu d'une population normale ; la variance est inconnue. Le test de Student est alors exact à toute taille. Une taille faible ne permet pas, à elle seule, de supposer la normalité.}
\item \question{Formuler les hypothèses pour les deux démarches.}
\reponse{M. A : $H_0:\mu=28$ contre $H_1:\mu\neq28$. Mme B : $H_0:\mu\geq28$ contre $H_1:\mu<28$, calibré à la frontière $\mu=28$.}
\item \question{Donner la statistique et sa loi sous l'égalité de référence.}
\reponse{$T=\frac{\overline X-28}{\frac{S}{\sqrt{10}}}\sim St(9)$ pour $\mu=28$, avec $S^2=\frac{\sum(X_i-\overline X)^2}{9}$.}
\item \question{Donner les régions critiques au seuil de $5\%$.}
\reponse{M. A rejette si $|t_{\mathrm{obs}}|>t_{9;0{,}975}\simeq2{,}26216$. Mme B rejette si $t_{\mathrm{obs}}<-t_{9;0{,}95}\simeq-1{,}83311$.}
\item \question{Calculer les statistiques observées.}
\reponse{$\bar x=24{,}5$, $s^2=29{,}83333$, $s\simeq5{,}46199$ et $t_{\mathrm{obs}}\simeq-2{,}02636$.}
\item \question{Conclure dans les deux cas, en calculant les $p$-valeurs. Un non-rejet prouve-t-il l'égalité ? Un rejet établit-il une cause ?}
\reponse{La $p$-valeur bilatérale est $0{,}07336$ : M. A ne rejette pas l'égalité à $5\%$, sans la prouver. La $p$-valeur à gauche est $0{,}03668$ : Mme B rejette $H_0$ et les données étayent une moyenne inférieure à 28 sous le modèle.
Le choix du sens du test après lecture des données invaliderait cette comparaison des risques annoncés. Ces résultats n'établissent pas à eux seuls une cause de l'écart salarial.}
\end{enumerate}
}
